\section{Which Teaching Tendencies Distinguish Models?}
\label{sec:differences}
\label{sec:results}

Some teaching behaviors show recurring differences between models; others
show differences only in certain settings, shared responses, or inconclusive
results (RQ1; Table~\ref{tab:finding-types}). We report the evidence for each
behavior separately.

\begin{table}[t]
\centering\small
\begin{tabular}{p{0.22\linewidth}p{0.34\linewidth}p{0.35\linewidth}}
\toprule
Finding type & Representative evidence & Interpretation\\
\midrule
Recurring helping differences & Answer revelation and reasoning invitations; complementary directness screen & Some helping tendencies distinguish models under the tested conditions\\
More limited recurrence & Response complexity, expressed confidence, and within-task safety patterns & Recurrence has different scopes across behaviors and measures\\
Shared responses & Identical default affiliation-inventory choices & These fixed-choice items do not distinguish the tested models\\
Unresolved consistency & Mixed revision/abstention results; failed candidate checks & Insufficient support for recurrence does not establish absence of differences\\
\bottomrule
\end{tabular}
\caption{Types of findings in Section~\ref{sec:differences}. Shared responses and
unresolved consistency are distinguished within the third subsection. The
full evidence map, including subsequent persistence tests, is in Appendix
Table~\ref{tab:behavior-map}.}
\label{tab:finding-types}
\end{table}

\subsection{Helping Choices Show Recurring Model Differences}

Models differ in how directly they help and how often they invite student
reasoning. On the 32 source problems under the common tutoring instruction, GLM reveals an answer in
22.3\% of responses and invites reasoning in 96.5\%; DeepSeek and Doubao reveal
answers in every response, with reasoning invitations below 5\%. M3 falls between these models on both actions (Appendix Table~\ref{tab:default-rates}). These rates describe usual behavior under the same role and inputs; they do
not rank teaching quality or assign models to fixed types. Giving an answer and inviting reasoning are separately coded actions:
an invitation leaves a content-related reasoning step for the student rather
than asking a question immediately answered by the tutor. A reply may contain
both actions without establishing their order or a multi-turn teaching process.

Ratings of assistance directness provide complementary evidence. The scale
runs from hints or questions to complete solutions, and model differences
recur across tasks (ICC 0.629). To pass the strongest screening rule, a
dimension must be measured reliably, distinguish models across tasks, and
receive further support from teaching actions, instruction changes, or
diagnosis outcomes. Only directness passes the full rule. Its changes under
instructions agree with independently inferred teaching actions, and the
direction remains the same when each coder is excluded
(Appendix~\ref{app:behavior-evidence}). The directness ratings and binary action
labels describe related but different aspects of helping. Section~\ref{sec:stability}
tests whether the action labels support prediction on new inputs.

\subsection{Other Differences Have More Limited Support}

Models also differ across tasks in the amount of information and number of
steps they present (ICC 0.663). These response-complexity ratings pass the
cross-task consistency criteria, but lack the additional evidence needed to
pass the full screening rule. They measure what a response contains, not the
mental effort a student experiences.

Differences in expressed confidence recur across four context families in six
historical configurations (ICC 0.852). This
exploratory archive result concerns how confidently models report their
answers; it does not establish correctness or calibration. Revision and abstention measures show mixed consistency, so expressed
confidence alone does not establish a broader cautiousness trait. Safety measures also show recurring differences within some tasks, but do not
support one ranking of how permissive models are across tasks. These results
apply to the reported measures and settings; they have not received the same
new-response tests as the three action measures (Appendix~\ref{app:behavior-evidence}).

\subsection{Shared Responses and Uncertain Differences}

Some measures yield common responses. In the affiliation pilot, all five models choose the affiliative option on
every default questionnaire item. Those choices do not distinguish models. Their freely written responses do show recurring differences within this
pilot, but affiliation ratings closely overlap with ratings of surface warmth.
The shared questionnaire choices therefore do not mean that the models behave
identically in open responses, and the pilot does not establish a separate
general affiliation trait (Appendix~\ref{app:behavior-evidence}).

Other measures fall short for different reasons. Ratings of how actionable the next step is change with instructions but fail
the cross-task stability check. The rated elicitation scale also fails its full
screening rule. The later predictive success of a separate binary measure---
whether a reply invites reasoning---does not validate that scale. We retain
these failed checks without treating them as proof that models never differ. The next experiments focus on predicting revelation, elicitation, and
acknowledgement and testing how instructions change them. They do not validate
every tendency examined in the broader screen.
